

\documentclass{beamer}
\usetheme{Berlin}
\usepackage{amsmath}

\usepackage{array}
\usepackage{booktabs}

\mode<presentation>{}
\usetheme{Warsaw}
\begin{document}


\frame{\titlepage} 
\section{Introduction}
\frame{
\frametitle{ImageNet Classification}
\begin{itemize}
\item ImageNet Classification with Deep Convolutional Neural Networks  
\end{itemize}
}


\begin{columns}
  \column{0.9\textwidth}
  \begin{block}{Input data}
    \begin{itemize}
    \item ILSVRC-2012 Challenge: 1.2 mill bilder fra ImageNet
      \item skalert ned til 256x256 pixeler
    \item Tensor med rank 3
      \begin{displaymath}
        X=X_{i,j,rgb} 
      \end{displaymath}
      \end{itemize}
\end{block} 
\end{columns}

\frame{
\frametitle{\O kning av datamenged for å unng\aa \ overfitting}
\section{Data augumentation}
Med  60 mill parametre og 1.2 mill bilder s\aa \space vil man f\aa \space problemer
med overfitting. Dette mitigeres ved å \o ke datamenged p\aa \space 2 m\aa ter. 

\begin{columns}
  \column{0.9\textwidth}
\begin{block}{Data augumentation}
\begin{itemize}
\item{Label preserving transformations}
  \begin{itemize}
  \item{Forkjellig utsnitt  224x224 }   
   \item Data \o kes med en faktor 2048
  \end{itemize}
  \item{PCA p\aa \space RGB fargeverdien (se neste slide) }
\end{itemize}
\end{block}
\end{columns}
}



\frame{
\frametitle{Tilfedlig bilde utsnitt}
\begin{columns}
  \column{0.9\textwidth}
  \begin{block}{Utsnitt og refleksjon}
    \begin{itemize}
    \item {Fra rammen på 256x256 tas tilfeldig utsnitt a 224x244 piksler}
    \item{Horisontal refleksjon av hvert bilde legges til data settet}
    \item{Antall bilder \o kes ned en faktor 2048}
    \item{Sterk korrelasjon, men hindrer overfitting}
\end{itemize}    
  \end{block}
\end{columns}
}


\frame{
\frametitle{ PCA }
\begin{columns}
  \column{0.9\textwidth}
  \begin{block}{PCA}
    \begin{itemize}
    \item Farge pixel (RGB)  verdien $I_{xy} = [ I_{rx}^R, I_{rx}^G, I_{rx}^B]$ endres ved å legge til st\o y langs egenvektorene til
    correlasjons matrisen $\mathbf{A}$ for RGB verdier.
    $$[\mathbf{p}_{1}, \mathbf{p}_{2}, \mathbf{p}_{3} ][\alpha_1 \lambda_1, \alpha_2 \lambda_2, \alpha_3 \lambda_3 ]^T$$    
    Egenvektoren er ortogonale (siden A er symetrisk).

  \item For hvert treningspass s\aa \space legges en st\o y verdi til langs egenvektorene p\aa \space hvert bilde.
    \end{itemize}
  \end{block}
\end{columns}
}

\section{Analyse}

\frame{
\frametitle{Analyse}
\begin{columns}
  \column{0.9\textwidth}
  \begin{block}{Utsnitt og refleksjon}
    \begin{itemize}
    \item {Hardware: 2xNivida GTX 580}
    \item {Deler av analyse kjeden g\aa r paralellt og uavhengig av hverandre på 2 GPU}
    \item {Trade off mellom turnover og kompleksitet} 
\end{itemize}    
  \end{block}
 \begin{block}{Simulering av nevroner}
    ReLU: \[f(x) = max(0,x)\] vs \[f(x) = 1/(1 + exp^-1)\]
    Trade off mellom turnover og n\o yaktighet. ReLU er 6 ganger rasker \aa \space trene   
  \end{block}
  
\end{columns}
}


\section{Main Results}
\frame{
\frametitle{Oppsummering}
\begin{columns}
  \column{0.9\textwidth}
  \begin{tabular}{ l  l  l  l }
  \toprule			
  & \multicolumn{2}{c}{Forbedring \%  } & \\
  Metode & top 1 & top 5 & Kommentar\\
  2XGPU &1.7  & 1.2  & vs halve kompleksiteten på 1 GPU\\
  Overlapping pooling  & 0.4 & 0.3 &  s=2, x=3  \\
  \bottomrule  
 \end{tabular}  
 \end{columns}
 }




\end{document}

